% Apéndice: bibliografía comentada por el profesor
\section{Bibliografía comentada por el profesor}\label{cap:20}\label{ap:bibliografia}

\textbf{Material docente.} Lista de la pestaña ``Bibliografía'' del campus, en el orden del
campus, con los comentarios textuales del profesor (entre comillas) y las indicaciones que
dejó en los resúmenes de clase. Los PDF están en la carpeta
\file{Proyectos / Agrimensura / Sistemas de Coordenadas / Bibliografia} de Dropbox; entre corchetes,
el nombre del archivo.

\begin{itemize}
  \item \textbf{Axler}, \emph{Linear Algebra Done Right} [\file{Algebra done right.pdf}].
  Resumen de clase 4: ``el libro del que yo primero saqué el tema [formas lineales] y
  concluí que era importante es el de Axler, que está en inglés''.
  \item \textbf{Gerlach}, \emph{Dual of a Vector Space} [\file{DualOfAVectorSpace.pdf}].
  Complemento para formas lineales y dualidad (sección~\ref{sec:formas-lineales}).
  \item \textbf{Butler}, \emph{Facts About Eigenvalues} [\file{butler-eigenvalues.pdf}].
  \item \textbf{Winitzki}, \emph{Linear Algebra via Exterior Products}
  [\file{Algebra via exterior products.pdf}]. Producto cuña y determinante
  (secciones~\ref{sec:exterior} y~\ref{sec:nucleo}).
  \item \textbf{Strang}, \emph{Álgebra lineal y sus aplicaciones} [\file{Strang-algebra-p1.pdf},
  sólo parte 1].
  \item \textbf{Juan de Burgos}, \emph{Álgebra Lineal y Geometría Cartesiana}
  [\file{Algebra-burgos-p1.pdf}, sólo parte 1]. ``Clase 3. A partir de la página 382 del
  libro de Burgos hay una muy buena explicación de geometría afín, aunque está basada toda en
  expresiones algebraicas. Es bastante importante notar que este enfoque es mixto, no se basa
  directamente en transformaciones lineales pero tampoco se basa en puntos y rectas
  abstractas sino que trabaja en coordenadas cartesianas.'' Resumen de clase 3: ``me parece
  que el que mejor les va a venir es el de Burgos, aunque el escaneo es francamente malo''.
  Resumen de clase 4: ``en castellano el libro que trata lo de formas lineales es el de
  Burgos (el capítulo que dice `espacio dual')''.
  \item \textbf{Azócar}, \emph{El cuerpo de los complejos}
  [\file{El Cuerpo de los Complejos - Mario Azócar.pdf}]. Rotaciones y complejos
  (sección~\ref{sec:rot-complejos}).
  \item \textbf{Rosenfeld}, \emph{Proyección Estereográfica}
  [\file{B. A. Rosenfeld - Proyección Estereográfica ... Mir - 1977.pdf}].
  \item \textbf{Noble y Daniel}, \emph{Álgebra Lineal Aplicada}
  [\file{Noble-Daniel-Algebra-p1.pdf}, sólo parte 1].
  \item \textbf{Markushévich}, \emph{Curvas Maravillosas}
  [\file{Markushevich-Curvas_Maravillosas.pdf}].
  \item \textbf{Coxeter}, \emph{Geometría} [\file{Geometria_Coxeter_v1.pdf}]. ``El libro
  de Coxeter es un libro clásico de geometría. Clase 3. El enfoque que da en el capítulo de
  geometría afín se relaciona con la forma en la que explicamos nosotros en lo siguiente: no
  depende de la expresión en coordenadas para definir las transformaciones. Nuestro enfoque
  fue menos riguroso (sobre todo por cuestiones de tiempo), pero es importante ver que a
  partir de puntos y rectas sueltos en el plano se pueden definir las transformaciones
  afines, su composición, etc.''
  \item \textbf{Fortney}, \emph{A Visual Introduction to Differential Forms and Calculus}
  [\file{Fortney,_J_P_A_Visual_Introduction_to_Differential_Forms_and_Calculus.pdf}].
  Formas lineales dibujadas como familias de rectas y planos; la figura de la
  sección~\ref{sec:formas-lineales} sigue ese estilo.
  \item \textbf{Efimov}, \emph{Formas Cuadráticas y Matrices} [\file{Efimov - Formas Cuadraticas.pdf}].
  \item \textbf{Goloviná}, \emph{Álgebra Lineal y algunas de sus aplicaciones}
  [\file{Golovina-Algebra-p2.pdf}, sólo parte 2].
  \item \textbf{Richter-Gebert}, \emph{Perspectives on Projective Geometry}
  [\file{Jürgen Richter-Gebert ... Springer (2011).pdf}]. Homogéneas, incidencia y
  dualidad (bloque~\ref{sec:bloque-proyectivas}).
  \item \textbf{Calvetti y Somersalo}, \emph{Less is More Linear Algebra}
  [\file{Calvetti - Somersalo Less is More Linear Algebra.pdf}].
  \item \textbf{Hitchman}, \emph{Geometry with an introduction to cosmic topology}
  [\file{Michael_P_Hitchman_Geometry_with_an_Introduction_to_Cosmic_Topology.pdf}].
  \item \textbf{Sáenz}, \emph{Cálculo Vectorial} [\file{Calculo Vectorial (Jorge-Saenz) 1a edicion.pdf}].
  \item \textbf{Mahecha}, \emph{Mecánica Clásica Avanzada}
  [\file{Mecanica Clasica Avanzada -- Jorge Mahecha Gómez.pdf}]. Resumen de clase 1:
  ``todo el tema de rotaciones de tres dimensiones es medio complejo de encontrar así en los
  libros (y menos en castellano), pero en otro orden lo pueden encontrar en el libro
  \emph{Mecánica Clásica Avanzada}, Jorge Mahecha, capítulo 7''. También el artículo de
  Wikipedia sobre el grupo de rotación \(SO(3)\).
  \item \textbf{Miles Reid y Szendr\H{o}i}, \emph{Geometry and Topology}
  [\file{Miles_Reid-Sendroi_Balazs-Geometry_and_topology.pdf}]. ``Clase 2 y Clase 3.
  Los capítulos 1 y 2 del libro de Reid explican muy bien cómo se pueden formar todas las
  similaridades del plano con reflexiones, es muy bueno para entender la composición de
  transformaciones y transformaciones en el plano en general. Está en inglés pero es un libro
  relativamente autocontenido y fácil en cuanto al lenguaje (no en cuanto a la matemática).''
  \item \textbf{Behnke y otros}, \emph{Fundamentals of Mathematics, vol.~2: Geometry}
  [\file{Behnke_Fundamentals_V2_Geometry_p1.pdf}, sólo parte 1]. ``Clase 3. En este
  libro, planteado desde el punto de vista de la geometría, el capítulo 3 explica
  transformaciones afines desde las propiedades y transformaciones de rectas (algo similar a
  lo que hicimos en clase), y recién al final arma la representación con matrices. Es
  complicado porque está en inglés, pero por ahora lo dejo porque tiene gráficos que están
  buenos. Clase 2 y 3. El capítulo 2 también está muy bueno, a partir de reflexiones
  construye todas las transformaciones de similaridad (afinidades que no deforman ángulos).
  Dejo el libro completo, el archivo principal es la primera parte, una vez que lo abren si
  le cambian `p1' por `p2' en la URL tienen la parte 2.''
  \item \textbf{Polya}, \emph{Cómo plantear y resolver problemas}
  [\file{Polya G - Como Plantear Y Resolver Problemas.PDF}]. Resumen de clase 7: ``de
  pasada hablamos sobre estrategias para resolver problemas [\ldots] les dejo en la sección
  de bibliografía un libro muy muy bueno sobre el tema''.
  \item \textbf{Fraleigh y Beauregard}, \emph{A first course in linear algebra} (inglés)
  [\file{Fraleigh_Beauregard-A_first_course_in_linear_algebra.pdf}] y
  \emph{Álgebra Lineal} (español, carpeta en 4 partes).
  \item \textbf{Beezer}, \emph{A first course in linear algebra} (en la carpeta
  \texttt{Bibliografia} del proyecto).
  \item \textbf{Lay}, \emph{Álgebra lineal y sus aplicaciones}
  [\file{Lay - Algebra lineal y sus aplicaciones.pdf}]. ``Un libro muy popular en internet
  (aunque no me guste tanto). Este libro tiene el ejemplo de coordenadas homogéneas cerca del
  final del capítulo 2 (Aplicaciones a gráficos por computadora). Además, los primeros dos
  capítulos sirven como introducción, bastante enfocada a resolver operaciones, pero de muy
  fácil lectura. (No tiene el enfoque geométrico que le damos en la materia tan marcadamente,
  pero eso no lo tiene casi ningún libro).''
\end{itemize}

\paragraph{Artículos y guías citados en los resúmenes} (misma carpeta de Dropbox).
\begin{itemize}
  \item \textbf{Pamfilos}, \emph{Isometries}, \emph{Similarities}, \emph{Affinity},
  \emph{Projective Plane} (clase 3): ``tienen muy buenos gráficos para todos estos ejemplos
  que hicimos en el pizarrón. Fíjense que en este libro se plantea un tipo de afinidad
  parecida a la que usamos en clase pero que tiene una línea fija y como guía una familia de
  rectas paralelas que no son necesariamente perpendiculares a la línea fija. Con este tipo
  de transformación es mucho más fácil mostrar que siempre hay una afinidad para llevar un
  triángulo a otro triángulo''.
  \item \textbf{Gentile}, \emph{Complejos y Transformaciones} y \emph{Complejos y
  Generalización}; \textbf{Herrera}, \emph{Complejos y Transformaciones} (clase 1): ``dos
  excelentes artículos de sendos gigantes de la matemática argentina''. Pendientes de bajar
  del campus.
  \item \textbf{Amoretti}, \emph{Equipolencia} (clase 4): ``un trabajo muy viejo [\ldots]
  relativamente fácil de leer y creo que puede aportar a pensar las cosas desde otro
  ángulo''. Grassmann, \emph{Extension Theory} (Google Books): ``lo más importante son los
  primeros dos capítulos''.
  \item \textbf{Allardice}, \emph{The barycentric calculus of Möbius}; \textbf{Möbius},
  \emph{On the composition of line segments} (clase 6): ``plantean las mismas ideas de
  linealidad pero no a partir de poner coordenadas sino a partir de hacer operaciones
  geométricas con segmentos''.
  \item \textbf{Molina y Oktaç}, \emph{Concepciones de la transformación lineal en contexto
  geométrico}; \textbf{González, Rojas, Roa y Fuentes}, \emph{Un esquema de transformación
  lineal} (clase 7): ``hablan de distintas formas de generar intuición sobre esto [el cambio
  de base]''.
  \item \textbf{Clavijo}, \emph{Apunte Teoría Elipsoide} (guía para clase, CC-BY-SA):
  geometría del elipsoide y cálculos geodésicos; no citado en los resúmenes, útil para el
  capítulo~\ref{cap:11}.
\end{itemize}

\paragraph{Videos recomendados en los resúmenes.} Fórmula de Euler
(\texttt{youtu.be/tOG\_vDb03T0}); exponencial de matrices
(\texttt{youtu.be/O85OWBJ2ayo}, \texttt{youtu.be/9CBS5CAynBE}); serie de cuatro videos sobre
transformaciones afines a partir de transformaciones lineales
(\texttt{youtu.be/fWRm9dISpNk}); videos de clase del profesor sobre producto cuña y SVD
(embebidos en el campus).

